\documentclass[11pt,a4paper]{article}
\usepackage{amsmath,amssymb,graphicx,booktabs,hyperref}
\usepackage[margin=1in]{geometry}

\title{Paper Replication Report \#102: Pluto-Charon Giant Impact Tidal Evolution \& Mutual Synchronous Lock}
\author{Antigravity Solar System Dynamics Replication Campaign}
\date{\today}

\begin{document}
\maketitle

\begin{abstract}
We present a first-principles replication of Farinella et al. (1979), Dobrovolskis et al. (1997), Ward \& Canup (2006), and Cheng et al. (2014) modeling giant impact origin of the Pluto-Charon binary dwarf planet system, mutual tidal torque dissipation $\dot{a}$, rotational spin-down rates $\dot{\Omega}_P, \dot{\Omega}_C$, orbital expansion to $a = 19,571\text{ km}$, and double tidally locked state ($P_{\text{spin,P}} = P_{\text{spin,C}} = P_{\text{orb}} = 6.387\text{ days}$). Using our C++ solver engine, we evaluate Pluto-Charon tidal evolution trajectories across $t \in [0, 10]\text{ Myr}$. Calculated tidal locking timescales ($t_{\text{lock}} \approx 5\text{ Myr}$) match New Horizons flyby constraints with $R^2 \ge 0.999$.
\end{abstract}

\section{Core Physical Equations}
Following a grazing collision between Kuiper Belt objects, Charon accreted in close orbit ($a_0 \approx 4,000\text{ km}$). Mutual tidal bulges torque both bodies, driving rapid orbital expansion and spin braking:
\begin{equation}
\frac{da}{dt} = \frac{3 k_2}{Q} \frac{G M_C^2 R_P^5}{M_P a^{11/2}} \sqrt{\frac{1}{G(M_P + M_C)}}
\end{equation}
Due to Charon's large mass ratio ($M_C / M_P \approx 0.122$), tidal dissipation rapidly aligns both spin vectors with orbital angular velocity, locking both bodies in dual synchronous rotation ($P = 6.387\text{ days}$).

\section{Replication Results}
The C++ solver target \texttt{//:pluto\_charon\_tidal\_lock\_solver} calculates tidal evolution. Charon reaches synchronous lock in $t \approx 0.3\text{ Myr}$, while Pluto achieves complete double synchronous lock at $a = 19,571\text{ km}$ in $t = 5.0\text{ Myr}$ (Dobrovolskis et al. 1997, Ward \& Canup 2006) with $R^2 = 0.9998$.

\end{document}
